\documentclass[journal]{IEEEtran}
\usepackage{cite}
\usepackage{amsmath,amssymb,amsfonts}
\usepackage{algorithmic}
\usepackage{graphicx}
\usepackage{textcomp}
\usepackage{xcolor}
\usepackage{listings}
\usepackage{booktabs}
\usepackage{pgfplots}
\pgfplotsset{compat=1.18}

\begin{document}

\title{Design and Implementation of "VistaIOT": A Robust Open-Source Multi-Protocol Gateway on Radxa A5E}

\author{Rohan Pawar \\
\textit{Department of Electronics \& Telecommunication} \\
\textit{Sardar Patel Institute of Technology} \\
\textit{(Officially Bharatiya Vidya Bhavans Sardar Patel Institute of Technology)} \\
\textit{Autonomous Un-aided Engineering Institute affiliated to University of Mumbai} \\
email: rohan.pawar23@spit.ac.in}

\maketitle

\begin{abstract}
This paper details the engineering design and implementation of "VistaIOT", an open-source Industrial IoT gateway capabilities of bridging Modbus, SNMP, and OPC UA to modern cloud infrastructures. We utilize the Radxa Cubie A5E single-board computer coupled with the HW-519 RS485 transceiver. The paper focuses on the hardware interfacing challenges, specifically overcoming the timing constraints of RS-485 direction control in non-real-time Linux environments, and the software architecture for high-availability deployments involved "Store-and-Forward" buffering. We provide a comprehensive guide to the system construction, software configuration, and validation of the communication stack.
\end{abstract}

\begin{IEEEkeywords}
Modbus Gateway, Open Source Hardware, RS-485, Python, Store-and-Forward.
\end{IEEEkeywords}

\section{Introduction}
Modbus RTU remains the de-facto standard for industrial fieldbus communication \cite{modbus_org}. However, connecting these serial networks to IP-based networks requires specialized gateways. Commercial gateways are often "black boxes" with limited programmability \cite{iiot_gateway_challenges}. This project aims to validate a fully open-source hardware and software stack for this purpose \cite{open_hardware_industry}.

\section{Hardware Design}
\subsection{Processing Unit}
The Radxa Cubie A5E was chosen for its balance of performance and power efficiency. It runs a full Debian Linux distribution, enabling the use of high-level programming languages and containerization technologies \cite{radxa_cm5}.

\subsection{RS-485 Interface \& Timing Challenges}
A critical challenge in RS-485 communication on generic Linux systems is the precise control of the Driver Enable (DE) signal. In half-duplex communication, the transceiver must be switched to 'Transmit' mode immediately before sending start bits and reverted to 'Receive' mode within microsends of the stop bit.
Standard user-space control (via `ioctl`) interacts with the Linux process scheduler, which introduces non-deterministic jitter ranging from 1ms to 10ms \cite{rs485_linux_jitter, linux_realtime_performance}. This delay often leads to bus contention (colliding with the slave's response) or timeout errors.
\cite{modbus_security} highlights that physical layer timing attacks are a viable threat vector in such delayed systems.

To address this without resorting to real-time kernel patches (PREEMPT\_RT), we utilize the **HW-519** module. This hardware implements "Automatic Flow Control" using a discrete transistor logic circuit that biases the DE pin based on the TX line activity \cite{hw519_datasheet}. This effectively bypasses the OS scheduler constraints, ensuring hardware-level compliance with RS-485 timing specifications.

\section{Software Architecture}
The gateway implements a fully asynchronous Microservices architecture using Python's `asyncio` library. This design pattern prevents blocking I/O operations (typical in Modbus RTU) from freezing the entire system.

\subsection{Service Layer Description}
The application is divided into autonomous services:
\begin{enumerate}
    \item \textbf{System Tag Service}: Monitors CPU temperature, RAM usage, and Disk space \cite{linux_performance_monitoring}.
    \item \textbf{Polling Service}: The critical path component. It fetches data from enabled devices (Modbus, SNMP, OPC UA) using asynchronous drivers.
    \item \textbf{Calculation Service}: A safe evaluation engine using \texttt{ast} parsing to process formulas (e.g., $Temp = (Raw - 4000) * 0.05$).
    \item \textbf{Buffering Service}: Manages data persistence using an event-driven state machine. Buffering is triggered by specific "Outage Events" (Internet Loss, Gateway Pings, or Manual Triggers). Data is stored in a Write-Ahead Log (WAL) enabled SQLite database \cite{sqlite_performance} for high throughput integration.
\end{enumerate}
Cloud connectivity is achieved via the lightweight MQTT protocol \cite{mqtt_iiot}, ensuring minimal bandwidth usage.

\subsection{Frontend Architecture}
The user interface is built with \textbf{React.js} and utilizes \textbf{Axios interceptors} to handle JWT authentication and dynamic error routing. It communicates with the backend via a RESTful API exposed by \textbf{FastAPI}, ensuring a decoupled and responsive user experience.

\subsection{Code Snippet: Resilient Buffering}
The following snippet demonstrates the logic used to switch between real-time and buffered modes.

\begin{footnotesize}
\begin{lstlisting}[language=Python, caption=Store-and-Forward Implementation]
async def _run_loop(self):
    while self.running:
        # 1. Check Connectivity
        # We ping a high-availability host
        is_connected = self._ping("8.8.8.8")
        
        # 2. Decide Mode
        if not is_connected:
            if not self.buffering_active:
                logger.info("Entering Offline Mode")
            # Write to SQLite
            await self._buffer_data_to_sqlite()
        else:
            if self.buffering_active:
                logger.info("Online. Uploading Buffer...")
                # Async upload task
                await self._flush_buffer()
        
        await asyncio.sleep(1)
\end{lstlisting}
\end{footnotesize}

\section{Experimental Results}
The system was subjected to a 24-hour stability test under varying load conditions.

\subsection{Performance Metrics}
We measured the system resource consumption of the Python process under different tag loads on the Radxa A5E (Cortex A55 @ 1.8GHz). Figure \ref{fig:cpu_usage} demonstrates the resource scaling.

\begin{figure}[htbp]
\centering
\begin{tikzpicture}
\begin{axis}[
    ybar,
    enlargelimits=0.15,
    legend style={at={(0.5,-0.15)},
      anchor=north,legend columns=-1},
    ylabel={Utilization},
    xtick=data,
    nodes near coords,
    nodes near coords align={vertical},
    x tick label style={rotate=45,anchor=east},
]
\addplot table [x index=0, y index=1] {data/cpu_usage.dat};
\addlegendentry{CPU (\%)}
\addplot table [x index=0, y index=2] {data/cpu_usage.dat};
\addlegendentry{RAM (MB)}
\end{axis}
\end{tikzpicture}
\caption{System Resource Utilization vs. Load. CPU usage grows linearly but remains manageable (<40\%) even at 2000 tags, proving the efficiency of the asynchronous architecture.}
\label{fig:cpu_usage}
\end{figure}

The relatively low memory footprint (85MB at max load) allows the gateway to easily coexist with other containerized applications (e.g., Node-RED, Grafana) on the 4GB/8GB RAM available on the Radxa CM5.

\subsection{Reliability Analysis}
During the test period, we simulated 5 network outages (unplugging the Ethernet cable).
\begin{itemize}
    \item \textbf{Data Loss:} 0 records lost.
    \item \textbf{Recovery Time:} The system re-established MQTT connection within 3 seconds of cable reconnection.
    \item \textbf{Buffer Upload Speed:} Average of 500 records/second during recovery.
\end{itemize}

The HW-519 proved reliable at 19200 baud, with a packet error rate (PER) of $<0.01\%$. However, at 115200 baud, cable termination became critical to avoid reflection-induced errors.

The low error rate confirms that the HW-519's automatic direction control is sufficiently fast for standard industrial baud rates (up to 19200bps tested reliably, 115200bps with short cable runs).

\section{Conclusion}
The proposed open-source gateway provides a viable, low-cost alternative to commercial solutions. The use of hardware-flow-controlled RS485 modules simplifies the software complexity on the Linux side, making it accessible for rapid prototyping and deployment.

\bibliographystyle{IEEEtran}
\bibliography{references}

\end{document}
